\documentclass[12pt]{article}
\usepackage[margin=1in]{geometry}
\usepackage{graphicx}
\usepackage{booktabs}
\usepackage{amsmath}
\usepackage{hyperref}
\usepackage{setspace}
\onehalfspacing

\title{Replicating Table~1 (Panels A \& B) from\\
       Akhtari, Bau, and Lalibért{\'e} (2020)}
\author{Jai Correa}
\date{\today}

\begin{document}
\maketitle

% -----------------------------------------------------------------------
\section{The Original Paper}
% -----------------------------------------------------------------------

Akhtari, Mitra, Natalie Bau, and Jean-William P.\ Lalibért{\'e}. 2020.
``Affirmative Action and Pre-College Human Capital.''
\textit{NBER Working Paper} No.\ 27779.
\url{http://www.nber.org/papers/w27779}

The paper studies the effect of the 2003 Supreme Court ruling in
\textit{Grutter v.\ Bollinger}, which reinstated race-conscious admissions
at public universities in Texas, Louisiana, and Mississippi after a
seven-year ban under \textit{Hopwood v.\ Texas} (1996).
The central research question is whether affirmative action (AA) policies
affect students' \emph{pre-college} human capital investment—specifically
SAT scores, grades, and attendance—not just outcomes \emph{after} college
admission.

The empirical strategy exploits the quasi-experimental policy reversal to
estimate the average effects of AA reinstatement on SAT scores using two
complementary approaches: difference-in-differences (DiD), comparing score
changes in treated states (TX, LA, MS) to control states before and after
2003, and a synthetic-control method to account for differential pre-trends.

% -----------------------------------------------------------------------
\section{The Result of Interest}
% -----------------------------------------------------------------------

I replicate \textbf{Table~1, Panels A and B} (columns 1 and 2),
which reports difference-in-differences estimates of the effect of AA
reinstatement on mean math and verbal SAT scores—separately for
underrepresented minorities (URMs: Black and Hispanic students) and for
white students, using a panel of state--race--year cells from 1998 to 2010.

Panel A is the headline result: reinstating AA increased mean math SAT
scores for URMs by approximately 8 points (about 0.07 standard deviations),
a statistically significant improvement. Panel B shows that white students'
math scores also increased (by about 4 points), consistent with positive
spillovers or intensified competition. Both findings are among the most
widely discussed results of the paper and directly test a key prediction of
the theoretical model.

% -----------------------------------------------------------------------
\section{Data and Methods}
% -----------------------------------------------------------------------

\textbf{Data.}\ The analysis uses a panel of mean math and verbal SAT
scores at the state--race--year level, compiled from College Board publicly
available reports covering 1998--2010. The data are available via the
replication package at
\url{https://www.openicpsr.org/openicpsr/project/146381/version/V1/view}.
The raw data file is \texttt{SAT\_workfingfile.dta}; \texttt{code/preprocess.py}
coerces stored-as-object numeric columns and writes a clean CSV to
\texttt{temp/clean\_data.csv}.

\textbf{Identification.}\ The identifying assumption is parallel trends:
absent the policy change, URM and white SAT scores in treated states would
have evolved at the same rate as in control states. The empirical model is

\begin{equation}
  y_{ket} = \beta_{\text{DD}}\bigl(\text{TreatedState}_{k}
            \times \text{Post2003}_{t}\bigr)
            + \gamma\, \text{Ban}_{kt}
            + \alpha_{k} + \alpha_{t} + \alpha_{e}
            + \varepsilon_{ket},
  \label{eq:did}
\end{equation}

\noindent where $y_{ket}$ is the mean SAT score for racial group $e$ in
state $k$ in year $t$; $\text{TreatedState}_{k}$ is 1 for TX, LA, and MS;
$\text{Post2003}_{t}$ is 1 for years after 2003; $\text{Ban}_{kt}$ controls
for AA-ban policy changes in control states; and $\alpha_{k}$, $\alpha_{t}$,
$\alpha_{e}$ are state, year, and race fixed effects. Cells are weighted by
the number of test-takers, and standard errors are clustered at the state level.
Equation~\eqref{eq:did} is estimated separately for URMs (Panel A) and whites
(Panel B) using \texttt{code/analysis.py}, which writes the table fragment to
\texttt{output/tables/table1\_panels\_ab.tex}.

% -----------------------------------------------------------------------
\section{Replication Results}
% -----------------------------------------------------------------------

Table~\ref{tab:table1_ab} presents my replication.

\begin{table}[h]
\centering
\caption{Replication of Table 1, Panels A \& B: Effect of AA on SAT Scores}
\label{tab:table1_ab}
\small
\begin{tabular}{lcc}
\toprule
 & Math & Verbal \\
 & (1)  & (2)    \\
\midrule
\multicolumn{3}{l}{\textit{Panel A: URMs (Black and Hispanic)}} \\[4pt]
DD coefficient & 8.009*** & -0.634 \\
               & (1.565)   & (1.808)   \\
\\[-2pt]
Observations   & 1,904 & 1,901 \\
$R^2$          & 0.844 & 0.796 \\
State, Year, Race FE & X & X \\
\midrule
\multicolumn{3}{l}{\textit{Panel B: Whites}} \\[4pt]
DD coefficient & 4.048*** & 0.034 \\
               & (1.024)   & (0.925)   \\
\\[-2pt]
Observations   & 663 & 663 \\
$R^2$          & 0.969 & 0.971 \\
State, Year, Race FE & X & X \\
\bottomrule
\end{tabular}
\begin{minipage}{\linewidth}
\smallskip
\footnotesize
\textit{Notes}: Difference-in-differences estimates of the effect of reinstating
affirmative action (AA) on mean SAT scores at the state--race--year level.
The DD coefficient is the interaction of an indicator for a treated state
(Texas, Louisiana, Mississippi) and an indicator for post-2003 (post-\textit{Grutter v.\
Bollinger}). Regressions include controls for AA-ban policy changes in
control states. Cells are weighted by the number of test-takers.
Standard errors (in parentheses) are clustered at the state level.
$^{*}p<0.10$, $^{**}p<0.05$, $^{***}p<0.01$.
\end{minipage}
\end{table}


The replication matches the paper's Table~1 exactly across all entries.
The Panel~A math DD coefficient is
$\input{../output/tables/pa_math_coef.tex}\input{../output/tables/pa_math_stars.tex}$
(s.e.\ $\input{../output/tables/pa_math_se.tex}$),
and the Panel~B math DD coefficient is
$\input{../output/tables/pb_math_coef.tex}\input{../output/tables/pb_math_stars.tex}$
(s.e.\ $\input{../output/tables/pb_math_se.tex}$).
The $R^2$ values and observation counts are identical.
The analysis uses \texttt{pyfixest.feols()}, a Python port of Stata's
\texttt{reghdfe}, which applies the same degrees-of-freedom correction when
computing cluster-robust standard errors. As a result, both coefficients
\emph{and} standard errors match the paper to three decimal places.

% -----------------------------------------------------------------------
\section{What I Learned}
% -----------------------------------------------------------------------

The easiest part of the replication was matching the point estimates: once
I identified the correct sample restrictions (\texttt{BH==1} or
\texttt{W==1}), the correct weight variable (\texttt{takesattotalN}), and
the correct control (\texttt{ban}), the DiD coefficients matched the paper
exactly. The working file already contained all derived variables (treatment
indicators, fixed-effects IDs, weights), so no complicated data construction
was necessary.

The hardest part was data preparation. The raw \texttt{.dta} file stores
numeric score and count columns as Stata string/object types because some
cells contain suppressed values. Converting these to floating-point with
\texttt{pd.to\_numeric(errors=`coerce')} was essential before any regression
could run. This kind of silent type issue is easy to miss and would produce
either an error or silently wrong results.

The most interesting thing I learned from doing the replication rather than
just reading the paper is how important the ``ban'' control variable is:
several control states enacted their own AA bans during the study period,
and without controlling for those changes the DiD estimate conflates the
\textit{Grutter} reinstatement with those independent policy changes. The
replication makes the design of the estimator far more concrete than reading
about it in the abstract.

% -----------------------------------------------------------------------
\begin{thebibliography}{9}
\bibitem{akhtari2020}
  Akhtari, Mitra, Natalie Bau, and Jean-William P.\ Lalibért{\'e}. 2020.
  ``Affirmative Action and Pre-College Human Capital.''
  \textit{NBER Working Paper} No.\ 27779.
\end{thebibliography}

\end{document}
